% =====================================================================
%  PLA DE TREBALL (PT) -- Concurs LE-403/729 (UPC) -- SKELETON
%  Structure drafted 15/11/2026 from the analysis of PT_Royano, PT_Joan and
%  Projecte_Professor_Agregat, adapted to the HAP (3 research lines).
%  Orange [Note: ...] = what to write and where it comes from.
%  Red [TODO: ...]    = information still to be provided.
%  Target length: 16--20 pages (Royano 17, Joan 25). American English.
%  Personal data: ONLY name and surname (base 3.5).
% =====================================================================
\documentclass[11pt]{article}
\usepackage{lector}
\newcommand{\sectionbreak}{\clearpage}
%\finaltrue
\setupheader{PT}

\title{\textbf{Pla de Treball}\\[2pt]\large Work Plan}
\author{\Large \CandName}
\date{}

\begin{document}
\maketitle
\thispagestyle{plain}
\vspace{0.5cm}
\begin{center}\postdata\end{center}

\vspace{0.8cm}
\noindent\textbf{Summary.}
\NOTE{One paragraph (8--10 lines). State how the plan answers the three requirements of the
call: (i) teaching in the thermofluid subjects of the EEBE and the Master's Degree in
Mechanical Technologies; (ii) research in the three lines of the department; (iii)
supervision of doctoral theses. Name the three research lines exactly as in HAP 3.3 and the
subjects developed in Part 2. End with the six-year horizon (2027--2032) and the Year-3
evaluation.}

\newpage
\tableofcontents

% =====================================================================
\section{Research plan}

\subsection{Background and motivation}
\NOTE{1--1.5 pages. Royano's ``Antecedentes y justificaci\'on'' model: start from what the
PhD established and which open questions it left, then the engineering motivation. Reuse
the narrative of HAP 3.3 but with references to the literature.}
\begin{itemize}
  \item Engineering motivation: supercritical fluids in sCO$_2$ power and refrigeration
  cycles, regenerative cooling of rocket engines, thermal management, chemical processing;
  hypersonic aerothermal loads for space exploration. \TODO{3--5 references}
  \item What my PhD and postdoctoral work established: DNS methodology for transcritical
  wall turbulence [J5, J6], flow regimes and secondary flows in ducts [J2, J4, R1], the
  GPU-resident solver RHEA-X [J1].
  \item Open questions that define this plan:
  \begin{itemize}
    \item Which mechanisms identified in canonical flows govern heat transfer and pressure
    loss in engineering configurations?
    \item Which flow processes precede extreme wall-heat-flux events, and how early can
    they be anticipated?
    \item How can high-fidelity solvers sample rare events at exascale with real-fluid
    thermodynamics?
  \end{itemize}
  \item Link to the department: SCRAMBLE (ERC), GReCEF, the three lines of the call.
\end{itemize}

\subsection{Research lines and objectives}
\NOTE{2--2.5 pages. Same three lines as HAP 3.3, each mapped to the profile line(s) of the
call. Royano's format works well: ``Current status'' (one short paragraph, with the
supporting papers) and ``Strategy 2027--2032'' (two objectives per line, numbered O1--O6).
Keep objectives measurable.}

\subsubsection*{Line 1 -- Physics of multiphysics turbulence in supercritical and
compressible flows (profile line a)}
\begin{itemize}
  \item \textbf{Current status.} \TODO{PhD results; [J2]--[J7], [R1]}
  \item \textbf{O1.} \TODO{e.g.\ high-fidelity databases of heated/cooled channels and ducts
  across the pseudo-boiling line, including engineering-relevant geometries}
  \item \textbf{O2.} \TODO{e.g.\ mechanisms governing heat-transfer deterioration/enhancement
  and pressure loss; reference data for engineering models (regenerative cooling, sCO$_2$
  heat exchangers)}
\end{itemize}

\subsubsection*{Line 2 -- Causality and data-driven prediction of extreme events (profile
line b)}
\begin{itemize}
  \item \textbf{Current status.} \TODO{CTR 2026; [R2], [R3]; RES IM-2026-1-0034;
  collaborations (Chandra, Lozano-Dur\'an, Mart\'inez-S\'anchez, Yang)}
  \item \textbf{O3.} \TODO{e.g.\ a Lagrangian causality framework (transfer entropy,
  non-stationary formulations) to rank the precursors of extreme wall-heat-flux events}
  \item \textbf{O4.} \TODO{e.g.\ causality-informed predictors and surrogates; transfer to
  hypersonic boundary layers (Salento) and ocean turbulence/microplastics (UBC)}
\end{itemize}

\subsubsection*{Line 3 -- Distributed and GPU-accelerated high-fidelity CFD for exascale
computing (profile line c)}
\begin{itemize}
  \item \textbf{Current status.} \TODO{RHEA-X [J1]; surrogates [R4]; hackathons; 6 HPC
  allocations}
  \item \textbf{O5.} \TODO{e.g.\ scalable, performance-portable open-source solver on
  pre-exascale/exascale systems; directive-based vs.\ vendor-specific programming}
  \item \textbf{O6.} \TODO{e.g.\ efficient real-fluid property evaluation and in-situ
  Lagrangian tracking of large particle ensembles; reproducible benchmarks}
\end{itemize}

\subsubsection*{Long-term direction}
\NOTE{One paragraph: the three lines converge on high-speed aerothermodynamics for space
exploration and on industry-linked thermofluid design (as in the closing paragraph of HAP
3.3). This shows a research trajectory beyond the contract.}

\subsection{Work plan and timeline (2027--2032)}
\NOTE{2--3 pages. Royano's two-phase structure fits the Lector contract and its Year-3
evaluation: Phase 1 (Y1--Y3) consolidation and independence; Phase 2 (Y4--Y6) expansion,
PhD supervision and transfer. Joan's format for each WP: tasks, required resources,
expected outcomes. Add numbered milestones (H1, H2, ...) tied to papers, proposals and
theses, and a Gantt chart (reuse the pgfgantt chart of the current PT).}
\begin{itemize}
  \item \textbf{WP0 -- Resources, group building and management (Y1--Y6).} \TODO{}
  \item \textbf{WP1 -- Supercritical and compressible turbulence (O1, O2).} \TODO{tasks,
  resources, outcomes}
  \item \textbf{WP2 -- Lagrangian causality and extreme events (O3).} \TODO{}
  \item \textbf{WP3 -- Prediction and transfer to other flows (O4).} \TODO{UBC 2027,
  Salento 2028}
  \item \textbf{WP4 -- Exascale solver and real-fluid models (O5, O6).} \TODO{}
  \item \textbf{WP5 -- Mobility, collaborations, dissemination and transfer (Y1--Y6).}
  \TODO{}
\end{itemize}
\NOTE{Milestones table: e.g.\ H1 publication of [R1]/[R2] (Y1); H2 first own competitive
project as PI (Y1--Y2); H3 first PhD student (Y2); H4 open release of the exascale solver
version (Y3); ... Expected outcomes per WP in numbers (Q1 papers, datasets, software
releases, theses).}

\subsection{Resources, funding strategy and risk management}
\NOTE{1.5 pages. Combine Joan (list of target calls; computational resources; software)
and Royano (staged strategy + contingency plan). Be concrete and realistic.}
\begin{itemize}
  \item \textbf{Starting point.} SCRAMBLE (ERC StG) resources and GReCEF infrastructure.
  \TODO{project end date}
  \item \textbf{Competitive funding (as PI or co-PI).} AEI \emph{Generaci\'on de
  Conocimiento} (young-researcher modality); AGAUR; ERC Starting Grant (medium term);
  MSCA Doctoral Networks / Postdoctoral Fellowships (as host); Horizon Europe /
  EuroHPC; Mitacs and bilateral mobility programs. \TODO{check eligibility and dates}
  \item \textbf{Computing.} RES and EuroHPC calls; Digital Research Alliance of Canada via
  UBC; track record of 6 allocations ($>$10 M core and GPU hours).
  \item \textbf{Software and data.} RHEA-X (open source), thermophysical surrogate
  software [R4], CoolProp; FAIR data repositories (Zenodo, Figshare+, BLASTNet).
  \item \textbf{Industry and transfer.} \MAYBE{Smart Ocean Peniche}; \TODO{other companies
  or industrial doctorates (Doctorat Industrial)}
  \item \textbf{Contingency plan.} \NOTE{What happens if a proposal or a PhD grant fails:
  e.g.\ scope reduction to canonical cases, TFM students for auxiliary tasks, reuse of
  existing databases, smaller allocations.}
\end{itemize}

\subsection{PhD supervision and group building}
\NOTE{0.5--1 page. Proposed thesis topics (one per line or crossing lines), funding route
for each (FPI/FPU, AGAUR FI, Doctorat Industrial, project), co-supervision first, sole
supervision later; TFG/TFM as entry point. Mention the hackathon team-leading experience
carefully (no overclaiming).}
\begin{enumerate}
  \item \TODO{Thesis 1 -- Lines 1--2: mechanisms and prediction of extreme heat-transfer
  events in supercritical wall turbulence}
  \item \TODO{Thesis 2 -- Line 3: exascale real-fluid solver and Lagrangian sampling}
  \item \TODO{Thesis 3 (optional) -- industrial doctorate / sCO$_2$ components}
\end{enumerate}

\subsection{Mobility and collaborations}
\NOTE{0.5--1 page. Planned stays with dates and goals (UBC summer 2027, 3 months, Mitacs;
Universit\`a del Salento 2028, 6 months, teaching concentrated in one semester); ongoing
collaborations (CTR, Caltech/MIT, Penn State, Stanford/Ihme) and prospective ones (IST,
Smart Ocean). Explain what each brings to the plan, not only the names.}

\subsection{Open science, transfer and dissemination}
\NOTE{0.5 page. Target journals (JFM, Phys.\ Rev.\ Fluids, IJHMT, Phys.\ Fluids, CPC, JCP);
conferences (THMT-27, EFDC3 2027, APS DFD, ETMM, ECCOMAS); open-source releases; FAIR
datasets; outreach (Setmana de la Ci\`encia, UPC Open Days).}

% =====================================================================
\section{Teaching plan}

\subsection{Teaching background and personal vision}
\NOTE{0.5--1 page. Brief: subjects taught, hours, surveys, coordination support (details
are in the HAP). Then a short personal vision: linking fundamental physics with
engineering design; industrial cases (inverter cooling at Continental); simulation as a
learning tool; feedback-based learning (``predict, measure, explain'').}

\subsection{Context: the degrees and subjects of the call}
\NOTE{1 page. Joan's ``Contextualization'' kept short (do not copy the 15 pages of the
Agregat project on the EHEA). EEBE/EPSEVG degrees, where each subject sits in the study
plan, number of groups. Then a table: Subject -- Degree(s) -- ECTS -- My contribution.
Reuse the table of the current PT.}
\begin{itemize}
  \item Thermodynamics and Heat Transfer (820020, 6 ECTS); Thermodynamics.
  \item Fluid Mechanics (820013).
  \item Fluid Engineering (340058, 6 ECTS, EPSEVG).
  \item Thermal Engineering; Power Plants and Renewable Energies; Thermofluid-Dynamic
  Generation.
  \item Computational Fluid Mechanics (Master's Degree in Mechanical Technologies).
\end{itemize}

\subsection{Developed subject: Thermodynamics and Heat Transfer (TTC)}
\NOTE{3--4 pages. Royano's five-part structure for each developed subject. Base it on the
official course guide of 820020 and on the TTC plan of the Agregat project (prerequisites,
objectives, methodology, assessment, resources, syllabus, schedule), but in your own words
and showing what you add.}
\subsubsection*{Learning objectives} \TODO{}
\subsubsection*{Contents and syllabus} \TODO{units and weeks}
\subsubsection*{Methodology} \TODO{active learning feasible with large groups; real-fluid
property tools (CoolProp) and computational examples; industrial cases}
\subsubsection*{Assessment} \TODO{as in the course guide; continuous assessment}
\subsubsection*{Schedule and student workload} \TODO{60 h on-site + 90 h autonomous}

\subsection{Developed subject: Fluid Engineering (EPSEVG)}
\NOTE{2--3 pages. Same five-part structure. This is the strongest teaching argument: you
will coordinate it and teach groups M65 and M66 in spring 2027. Reuse the detailed section
of the current PT (objectives, 6 units, methodology, workload 52.5/7.5/90, assessment
CP1/CP2/labs/problems, bibliography). Note: the course guide lists Catalan as the teaching
language.}
\subsubsection*{Learning objectives} \TODO{}
\subsubsection*{Contents and syllabus} \TODO{}
\subsubsection*{Methodology} \TODO{external aerodynamics cases (Formula Student, truck drag,
NACA~0012 database [R5]); turbomachinery}
\subsubsection*{Assessment} \TODO{}
\subsubsection*{Schedule and student workload} \TODO{}

\subsection{Master's subject: Computational Fluid Mechanics}
\NOTE{1 page. Proposal of contributions agreed with the subject coordinators: GPU laboratory
session (CUDA/OpenACC on the subject's Poisson/Jacobi example), data-driven and causal flow
analysis, course projects that lead to TFM.}

\subsection{Teaching innovation and training plan}
\NOTE{1 page. Joan's ``Training Plan Innovation'' + ``Methodology Innovation'': concrete,
feasible actions (computational notebooks in TTC/MF/Fluid Engineering; laboratory guides
built around ``predict, measure, explain''; responsible use of generative AI following UPC
guidelines; flipped-classroom material shared among lecturers). ICE training pathway
(Atenea M2.4, conflict management, active methodologies) and participation in
teaching-innovation projects.}

\subsection{Supervision of TFG, TFM and PhD theses}
\NOTE{0.5 page. Example topics per level, each feeding a research line.}

\subsection{Language plan}
\NOTE{0.5 page. Current A2; B1 course scheduled in 2026; C1 within the first two years
(SLT pathway or CIFALC); temporary exemption requested under CG/2025/06/29. Teaching
languages and how the EPSEVG Catalan requirement will be met.}

\subsection{Equality, inclusion and sustainability}
\NOTE{0.5 page. Joan includes this; keep it concrete: inclusive examples and groups, UPC
gender-equality plan, sustainability in TFG topics (energy transition).}

% =====================================================================
\section{Year-3 milestones and self-assessment}
\NOTE{0.5--1 page. Table of measurable indicators for the intermediate evaluation:
publications (Q1, first/corresponding author), funding obtained as PI, PhD students
recruited, teaching (surveys, coordination, innovation projects), Catalan C1, mobility
completed. Add Year-6 targets in a second column.}

% =====================================================================
\section*{References}
\addcontentsline{toc}{section}{References}
\NOTE{Numbered references cited in Part 1 and in the subject plans. Reuse and verify the 13
references of the current PT.}

\end{document}
